\section{\texorpdfstring{Proofs for Sections~\ref{sec:general-start}
and~\ref{sec:frontier-modes}}{Proofs for the start and the field}}\label{app:start-field}

\subsection{The start weights}

We prove Lemma~\ref{lem:start-weights} and
Proposition~\ref{prop:start-weights}. Throughout $1/d<p<1$, so $\sigma=p-r>0$
and $v=r/\sigma=u/(1-u)$. The refresh form in Proposition~\ref{prop:dpmf}(a) is
linear in $\mu$, and with $\mu=\delta_a$ it gives $p^{N-1}S^{(a)}_n(u)$. So, with
sums over $1\le m_b\le n_b$ for $b\in F$ and $M=\sum_bm_b$,
\begin{equation}\label{eq:startw-weights}
 \omega_a(n)=\E_w\Bigl[\frac{m_a}M\Bigr],\qquad
 w(m)\propto\frac{M!}{\prod_bm_b!}\,v^M\prod_{b\in F}\binom{n_b-1}{m_b-1}.
\end{equation}
So $\omega_a(n)$ is the mean fraction of the blocks of the refresh form
that carry state $a$.

\begin{proof}[Proof of Lemma~\ref{lem:start-weights}]
(1) Exchanging the letters $a$ and $b$ is a bijection from the words with counts
$n$ that start with $a$ to those that start with $b$. It keeps the counts, since
$n_a=n_b$, and the number of changes, so $S^{(a)}_n=S^{(b)}_n$. If $k$ divides
$N$, a balanced $n$ has equal coordinates, and the $k$ weights are equal and sum
to $1$.

(2) We may take $N\ge N_0$ so large that $p>1/d$ and $\sigma\ge\frac12$. Let $n$
be balanced and $\nu=\lfloor N/k\rfloor$, so each $n_b$ is $\nu$ or $\nu+1$, and
$\nu\ge N/(2k)$.

\emph{Step 1. Reduction to two values.} By (1) the weights take a value $\omega_+$
where $n_b=\nu+1$ and a value $\omega_-$ where $n_b=\nu$. We show that
$0\le\omega_+-\omega_-\le\delta$ with $\delta=\frac2\nu(\E_wM+\E_wM^2)$. Since
the weights average to $1/k$, each is then within $\delta$ of $1/k$, and
$|\sum_a\mu_a\omega_a(n)-\mu(F)/k|\le\mu(F)\delta$.

\emph{Step 2. Pairing.} Let $n_a=\nu+1$ and $n_b=\nu$, and let $m'$ be $m$ with
$m_a$ and $m_b$ exchanged, with $w(m')=0$ if $m'$ is not admissible. All factors
of $w(m)$ other than $\binom{\nu}{m_a-1}\binom{\nu-1}{m_b-1}$ are symmetric in
$m_a$ and $m_b$. From $\binom\nu j/\binom{\nu-1}j=\nu/(\nu-j)$ we get
$w(m')=w(m)(\nu-m_a+1)/(\nu-m_b+1)$ when $m_a\le\nu$, and $w(m')=0$ when
$m_a=\nu+1$. In both cases $w(m)-w(m')=w(m)(m_a-m_b)/(\nu-m_b+1)$ when
$m_a>m_b$. The map $m\mapsto m'$ matches the admissible $m$ with $m_a<m_b$ to
the admissible $m$ with $m_b<m_a\le\nu$. Pairing $m$ with $m'$
in~\eqref{eq:startw-weights} gives, with $Z=\sum_mw(m)$,
\[
 \omega_a-\omega_b=\frac1Z\sum_{m_a>m_b}\bigl(w(m)-w(m')\bigr)\frac{m_a-m_b}M
 =\frac1Z\sum_{m_a>m_b}w(m)\frac{(m_a-m_b)^2}{M(\nu-m_b+1)}\ge0.
\]
Here $(m_a-m_b)^2/M\le M$. If $m_b\le\nu/2$ the last factor is at most $2/\nu$. If
$m_b>\nu/2$, then $M>\nu/2$ and the term is at most $w(m)M<w(m)\,2M^2/\nu$. This
proves the bound $\delta$.

\emph{Step 3. The moments.} It remains to bound $\E_wM$ and $\E_wM^2$. Since
$(m-1)!\binom{n_b-1}{m-1}/n_b^{m-1}=\prod_{j<m}(1-j/n_b)_+$, we have
$w(m)\propto c(m)A(m)$ for all $m$ with every $m_b\ge1$, where
\[
 c(m)=M!\prod_b\frac{(n_bv)^{m_b}}{m_b!(m_b-1)!},\qquad
 A(m)=\prod_b\prod_{j<m_b}\Bigl(1-\frac j{n_b}\Bigr)_+\in[0,1],
\]
and Lemma~\ref{lem:clipped-product} gives $1-A(m)\le\sum_bm_b^2/(2n_b)\le M^2/\nu$.
Write $M!=\int_0^\infty e^{-t}t^M\,dt$. Since
$\sum_{m\ge1}y^m/(m!(m-1)!)=y\Phi(y)$, this gives
$\E_c[(M+1)(M+2)]=\langle t^2\rangle$, where $\langle\cdot\rangle$ is the mean
under the density proportional to $e^{-t}t^k\prod_b\Phi(n_bvt)$, as in the
proof of Theorem~\ref{thm:uniform-simplex-majorant}. The integration by parts
there, with $t^{k+2}$ in place of $t^{k+1}$, gives
$\langle t^2\rangle=(k+2)\langle t\rangle+\sum_b\langle t\psi(n_bvt)\rangle$. Put
$c_1=\sqrt v\sum_b\sqrt{n_b}$. That proof gives
$\langle t\rangle\le2(k+1)+c_1^2$, and $\psi(x)\le\sqrt x$ and
$\langle t^{3/2}\rangle\le\langle t^2\rangle^{3/4}$ give
$X=\langle t^2\rangle\le(k+2)(2k+2+c_1^2)+c_1X^{3/4}$. If $c_1X^{3/4}\ge X/2$,
then $X\le16c_1^4$, and otherwise $X\le2(k+2)(2k+2+c_1^2)$. So
$X\le C_2(1+c_1^2)^2$, where $C_2$ depends only on $k$. Here
$c_1^2\le kNv=kT/\sigma\le2kT\le2kCL$, so
$\E_cM^2\le\E_c[(M+1)(M+2)]\le\nu/2$ for $N\ge N_0$. Since $0\le A\le1$ and
$\E_c[A]\ge1-\E_cM^2/\nu\ge\frac12$, and $M\le M^2$,
\[
 \E_wM\le\E_wM^2=\frac{\E_c[M^2A]}{\E_c[A]}
 \le\frac{\E_cM^2}{1-\E_cM^2/\nu}\le2C_2(1+2kT)^2.
\]
With $\nu\ge N/(2k)$ this gives $\delta\le C_1(1+T)^2/N$.
\end{proof}

\begin{proof}[Proof of Proposition~\ref{prop:start-weights}]
Take $N$ large, so that $0<u<1$ and $z=Nu\in[cL,2CL]$. Let $k=|F|$, put
$s=k-1$, $\Sigma=\sum_{b\in F}\sqrt{n_b}=A\sqrt N$ and
$\lambda=u(A/\sqrt{\alpha_a}-1)$. Note that
$A/\sqrt{\alpha_a}-1=\sum_{b\ne a}\sqrt{\alpha_b/\alpha_a}$ lies in
$[\sqrt\varepsilon,s/\sqrt\varepsilon]$. Let $j_1=\lfloor N/\sqrt L\rfloor$, so
that $j_1\le\varepsilon N/2$ for large $N$. We split the sum of
Lemma~\ref{lem:first-run} at $j_1$. The constants below depend only on $d$,
$k$, $\varepsilon$, $c$ and $C$.

\emph{Step 1. One term.} Let $1\le j\le j_1$, and put $n'=n-je_a$, $N'=N-j$ and
$\alpha'=n'/N'$. Then every $\alpha'_b\ge\varepsilon/2$,
$|\alpha'-\alpha|\le2\sqrt k\,j/N$ and $N'u\in[z/2,z]$. Since $N'\ge N/2$, for
large $N$ both $N'u$ and $z$ lie between $\frac c2\log N'$ and $4C\log N'$. So
Theorem~\ref{thm:simplex-interior}, with $\varepsilon/2$ in place of
$\varepsilon$ and with the constants $\frac c2$ and $4C$, applies at the totals
$N'$ and $N$ with errors $O(1/L)$. The difference of its exponents is
$(A(\alpha')^2-1)N'u-(A^2-1)z=u(\Sigma'^2-\Sigma^2)+ju$, where
$\Sigma'=\Sigma-\delta$ and
$\delta=\sqrt{n_a}-\sqrt{n_a-j}=j/(\sqrt{n_a}+\sqrt{n_a-j})$. So
$\delta\le j/\sqrt{n_a}$, and
\[
 0\le\delta-\frac j{2\sqrt{n_a}}
 =\frac{j\delta}{2\sqrt{n_a}\,(\sqrt{n_a}+\sqrt{n_a-j})}
 \le\frac{j^2}{2n_a^{3/2}}.
\]
Since $\Sigma\le\sqrt{kN}$ and $n_a\ge\varepsilon N$, this gives
$u(\Sigma'^2-\Sigma^2)=-2u\Sigma\delta+u\delta^2=-juA/\sqrt{\alpha_a}
+O(zj^2/N^2)$. Also $D(\alpha')/D(\alpha)=1+O(j/N)$, since $D$ is smooth and
positive on $\{\alpha_b\ge\varepsilon/2\}$, and the powers of $N'u$ and $N'$
change by the factor $(1-j/N)^{-s/2}=1+O(j/N)$. Hence there is $C_2$ with
\begin{equation}\label{eq:start-weights-term}
 S_{n-je_a}(u)=S_n(u)\,e^{-j\lambda+E_j},\qquad
 |E_j|\le C_2\Bigl(\frac1L+\frac{zj^2}{N^2}+\frac jN\Bigr).
\end{equation}

\emph{Step 2. The main sum.} For $j\le j_1$ we have $zj^2/N^2\le ju/\sqrt L$ and
$j/N\le ju/(cL)$. So for large $N$ we have $|E_j|\le1+j\lambda/2$, and then
$|e^{E_j}-1|\le e\,|E_j|e^{j\lambda/2}$. Put $\theta=\lambda/2$, so that
$u\sqrt\varepsilon/2\le\theta\le1$. Since $1-e^{-\theta}\ge\theta/2$, we have
$\sum_{j\ge1}j^me^{-j\theta}\le16\,\theta^{-m-1}$ for $m=0,1,2$. Hence
\[
 u\sum_{j\le j_1}(1-u)^{j-1}e^{-j\lambda}\bigl|e^{E_j}-1\bigr|
 \le eC_2u\sum_{j\ge1}e^{-j\theta}\Bigl(\frac1L+\frac{zj^2}{N^2}+\frac jN\Bigr)
 =O\Bigl(\frac u{L\theta}+\frac{uz}{N^2\theta^3}+\frac u{N\theta^2}\Bigr),
\]
which is $O(1/L)$. The geometric sum is
$u\sum_{j\ge1}(1-u)^{j-1}e^{-j\lambda}=u/(e^\lambda-1+u)$. Since $\lambda=O(u)$
and $\lambda+u=uA/\sqrt{\alpha_a}$, it equals $(\sqrt{\alpha_a}/A)(1+O(u))$. Its
terms with $j>j_1$ add up to at most
$u\sum_{j>j_1}e^{-j\lambda}\le(2u/\lambda)e^{-j_1\lambda}$, which is
$O(e^{-c\sqrt{\varepsilon L}/2})$. With~\eqref{eq:start-weights-term} we get
\[
 u\sum_{j\le j_1}(1-u)^{j-1}S_{n-je_a}(u)
 =S_n(u)\Bigl(\frac{\sqrt{\alpha_a}}A+O\Bigl(\frac1L\Bigr)\Bigr).
\]

\emph{Step 3. The other terms.} The map in Step~2 of the proof of the exact
assertions in Theorem~\ref{thm:simplex-crossing}, which adds a copy of $a$ just
before its first occurrence, shows that $S_m(u)\le S_{m+e_a}(u)$ for $u>0$
whenever $m_a\ge1$. So $S_{n-je_a}(u)\le S_{n-j_1e_a}(u)$ for $j_1<j<n_a$, and
by~\eqref{eq:start-weights-term} at $j=j_1$ these terms add up to at most
$S_n(u)e^{-j_1\lambda+E_{j_1}}$, since $u\sum_j(1-u)^{j-1}\le1$. Here
$zj_1^2/N^2\le z/L\le2C$ and $j_1/N\le1$, so $|E_{j_1}|\le C_2(1/L+2C+1)=O(1)$.
Also $j_1\lambda\ge(N/\sqrt L-1)u\sqrt\varepsilon\ge c\sqrt{\varepsilon L}/2$
for large $N$. So these terms are $O(S_n(u)e^{-c\sqrt{\varepsilon L}/2})
=O(S_n(u)/L)$. The term $j=n_a$ is at most
$u(1-u)^{n_a-1}S_{n''}(u)\le2ue^{-\alpha_az}S_{n''}(u)$, where $n''$ has the
$k-1$ coordinates $n_b$ with $b\ne a$, and total $N''=N(1-\alpha_a)$.
Theorem~\ref{thm:simplex-interior} gives
$S_n(u)\ge c_3z^{s/2}N^{-s}e^{(A^2-1)z}$ with $c_3>0$. If $k=2$, then
$S_{n''}=1$ and $A^2-1=2\sqrt{\alpha_1\alpha_2}\ge2\varepsilon$, so the term is
$O(S_n(u)z^{1/2}e^{-2\varepsilon z})$. If $k\ge3$, the fractions $n_b/N''$ are at
least $\varepsilon$ and $N''u\in[\varepsilon z,z]$. Since
$(\sum_{b\ne a}\sqrt{n_b})^2=N(A-\sqrt{\alpha_a})^2$,
Theorem~\ref{thm:simplex-interior} on $k-1$ states, with the constants
$\frac{\varepsilon c}2$ and $4C$, gives
$S_{n''}(u)\le C_3z^{(s-1)/2}N''^{1-s}e^{z[(A-\sqrt{\alpha_a})^2-1+\alpha_a]}$.
So the term is
$O(S_n(u)z^{1/2}e^{-z\sqrt{\alpha_a}(2A-\sqrt{\alpha_a})})
=O(S_n(u)z^{1/2}e^{-\varepsilon z})$. In both cases it is $O(S_n(u)/L)$.

Adding the three steps and dividing by $S_n(u)$ gives
$\omega_a(n)=S^{(a)}_n(u)/S_n(u)=\sqrt{\alpha_a}/A+O(1/L)$.
\end{proof}

\subsection{Gaussian confinement}

The next lemma gives the Gaussian profile of the ratio to a vertex near the
center of a face. We use it for Theorems~\ref{thm:start-maxima} and~\ref{thm:weakfield-phases}.

\begin{lemma}[Gaussian confinement]\label{lem:confinement}
Fix a support $S$ with $|S|=k\ge2$ and a compact set $\mathcal J$ of values of
$t$, put $T=L-\frac12\log L+t$, and let $u=r/p$ be the odds at this $T$. For a
positive count vector $n$ on $S$ put $y=\sqrt L\,(n/N-\frac1k\mathbf1)$, a
vector in the hyperplane $\{y\in\mathbb R^S:\sum_iy_i=0\}$. The terms $o(1)$
below tend to $0$ as $N\to\infty$.
\begin{enumerate}[(a)]
\item Uniformly for $t\in\mathcal J$ and $y$ in a compact subset of this
hyperplane,
\[
 \log S_n(u)=(k-1)(t-a_k)-\frac{k^2}4|y|^2+o(1).
\]
\item For every $R>0$, uniformly for $t\in\mathcal J$ and all $n$ with $|y|>R$,
\[
 \log S_n(u)\le(k-1)(t-a_k)-\frac{k^2}4R^2+o(1).
\]
\end{enumerate}
\end{lemma}

\begin{proof}
(a) Put $\alpha=n/N=\frac1k\mathbf1+y/\sqrt L$. For $|y|$ bounded,
$\alpha_i\ge1/(2k)$ for large $N$, and Taylor expansion of
$\sqrt{1/k+x}=k^{-1/2}+\frac{\sqrt k}2x-\frac{k^{3/2}}8x^2+O(x^3)$, with
$\sum_iy_i=0$, gives
\[
 \Bigl(\sum_i\sqrt{\alpha_i}\Bigr)^2=k-\frac{k^2}{4L}|y|^2+O(L^{-3/2}),
 \qquad D(\alpha)=D_k(1+o(1)).
\]
Insert these into Theorem~\ref{thm:simplex-interior}, with $\varepsilon=1/(2k)$,
$c=\frac12$, $C=2$ and $z=Nu=T/p$, where $z-T=O(L^2/N)$. Then
$(A^2-1)z=(k-1)z-\frac{k^2}4|y|^2+o(1)$ and
$z+\frac12\log z-L=t+\frac12\log(z/L)+o(1)=t+o(1)$. Since
$\log D_k=-(k-1)a_k$, this gives (a).

(b) For large $N$ we have $p>1/d$, since $T=O(L)$. Starting from $n$, move one unit from a coordinate $b$ to a coordinate $a$
with $n_b\ge n_a+2$, and repeat. The moves keep the support and end at a
balanced vector, and by Theorem~\ref{thm:simplex-balance} each move strictly
increases $S_n(u)$. A move changes $|y|^2$ by $\frac{2L}{N^2}(n_a-n_b+1)<0$ and
changes $y$ by a vector of length $\sqrt{2L}/N$. The balanced end has
$|y|=O(\sqrt L/N)$. So if $|y|>R$ at the start, the first vector on the path
with $|y|\le R$ has $|y|\ge R-\sqrt{2L}/N=R-o(1)$, and by (a), applied on the
compact set $\{|y|\le R\}$, its value of $\log S_n$ is at most
$(k-1)(t-a_k)-k^2R^2/4+o(1)$, uniformly. The starting vector has a smaller value.
\end{proof}

\subsection{Face maxima with a start}

Let $F$ be a face with $k\ge2$ states and put $s=k-1$. For $1\le j\le d$ and
$M\ge j$ let $\bar S_{M,j}(u)$ be the value of $S_m(u)$ at a balanced vector $m$
with $j$ positive coordinates and total $M$. So $\bar S_{N,k}=C_{N,k}/V_N$ and
$\bar S_{M,1}=1$. We use two facts. First, if $0<u\le1$, then
$S_m(u)\le\bar S_{M,j}(u)$ for every positive vector $m$ with $j$ coordinates
and total $M$. This is Theorem~\ref{thm:simplex-balance}, since
$S_m(u)=P_M(m)/V_M$ and $u\le1$ means $p\ge1/d$. Second, if $u>0$, then
$\bar S_{M,j}(u)\le\bar S_{M+1,j}(u)$, by Step~2 of the proof of the exact
assertions in Theorem~\ref{thm:simplex-crossing}.

\begin{lemma}[the start weights from above]\label{lem:start-upper}
Fix $0<c<C$. There is $C_4$ such that for large $N$, uniformly for
$cL\le T\le CL$,
\begin{equation}\label{eq:first-run-sum}
 u\sum_{1\le j\le N/2}(1-u)^{j-1}\bar S_{N-j,k}(u)
 \le\frac1k\Bigl(1+\frac{C_4}L\Bigr)\bar S_{N,k}(u),
\end{equation}
and $kS^{(a)}_n(u)\le(1+C_4/L)\,\bar S_{N,k}(u)$ for every positive count vector
$n$ on $F$ and every $a\in F$.
\end{lemma}

\begin{proof}
Take $N$ so large that $p>1/d$. Then $0<u<1$, and $z=Nu=T/p$ lies in
$[cL,2CL]$.

\emph{Step 1. The sum.} Let $1\le j\le N/2$ and $M=N-j$. Then $M\ge N/2$ and
$Mu=z(1-j/N)\in[z/2,z]$, so for large $N$ both $Mu$ and $z$ lie between
$\frac c2\log M$ and $4C\log M$. So~\eqref{eq:simplex-uniform-ratio} applies at
the totals $M$ and $N$, with the constants $\frac c2$ and $4C$, and its errors
are $O(1/L)$. Since $Mu-z=-ju$, it gives
\[
 \frac{\bar S_{M,k}(u)}{\bar S_{N,k}(u)}
 =\Bigl(1-\frac jN\Bigr)^{-s/2}e^{-sju}\Bigl(1+O\Bigl(\frac1L\Bigr)\Bigr).
\]
For $0\le x\le\frac12$ we have $\log(1-x)\ge-2x$, since the left side is concave
and the inequality holds at both ends. So $(1-j/N)^{-s/2}\le e^{sj/N}$. Put
$q=ku-s/N=ku(1-s/(kz))$, which is positive. With $(1-u)^{j-1}\le e^{-(j-1)u}$
and $\sum_{j\ge1}e^{-jq}=1/(e^q-1)\le1/q$, the left side
of~\eqref{eq:first-run-sum} is at most
\[
 \Bigl(1+O\Bigl(\frac1L\Bigr)\Bigr)\bar S_{N,k}(u)\,ue^u\sum_{j\ge1}e^{-jq}
 \le\Bigl(1+O\Bigl(\frac1L\Bigr)\Bigr)\bar S_{N,k}(u)\,\frac{ue^u}q .
\]
Here $ue^u/q=e^u/(k(1-s/(kz)))=(1+O(1/L))/k$, since $u=O(L/N)$ and $z\ge cL$.
This proves~\eqref{eq:first-run-sum}.

\emph{Step 2. The bound.} Apply the first fact to the vectors $n-je_a$ in
Lemma~\ref{lem:first-run}. They have $k$ positive coordinates for $j<n_a$ and
$k-1$ for $j=n_a$. With $(1-u)^{n_a-1}\le1$ and the second fact, this gives
\[
 S^{(a)}_n(u)\le u\sum_{1\le j\le N/2}(1-u)^{j-1}\bar S_{N-j,k}(u)
 +\bar S_{N,k}(u)\sum_{j>N/2}u(1-u)^{j-1}+u\,\bar S_{N,k-1}(u).
\]
The first term is bounded by~\eqref{eq:first-run-sum}. In the second,
$\sum_{j>N/2}u(1-u)^{j-1}\le(1-u)^{(N-1)/2}\le2e^{-z/2}\le2N^{-c/2}$. For the
third, Theorem~\ref{thm:simplex-window}(i), for $k$ and for $k-1$ when $k\ge3$,
gives $\bar S_{N,k-1}(u)/\bar S_{N,k}(u)=O(e^{-\Psi_N(T)})$. This also holds for
$k=2$, since $\bar S_{N,1}=1$. Here $ue^{-\Psi_N(T)}=zT^{-1/2}e^{-T}
=O(L^{1/2}N^{-c})$. So the second and third terms are
$O(\bar S_{N,k}(u)/L)$, which proves the bound.
\end{proof}

\begin{proof}[Proof of Theorem~\ref{thm:start-maxima}]
(1) Let $n$ have support $F$. By Proposition~\ref{prop:dpmf}(b) and
Lemma~\ref{lem:start-upper},
\[
 \frac{P^\mu_N(n)}{p^{N-1}}=\sum_{a\in F}\mu_aS^{(a)}_n(u)
 \le\mu(F)\max_{a\in F}S^{(a)}_n(u)
 \le\frac{\mu(F)}k\Bigl(1+\frac{C_4}L\Bigr)\bar S_{N,k}(u).
\]
By Theorem~\ref{thm:start-crossings}(2), a balanced bin of $F$ has
$P^\mu_N(n)/p^{N-1}=\frac{\mu(F)}k\bar S_{N,k}(u)(1+O(L^2/N))$. Since
$p^{N-1}\bar S_{N,k}(u)=d\,C_{N,k}$, this proves the first display. The second
follows from Theorem~\ref{thm:simplex-window}(iii), which gives
$\log\bar S_{N,k}(u)=(k-1)(t-a_k)+o(1)$.

(2) Fix $\mathcal J$, $R>0$ and $K\ge1$. Let $n$ have support $F$ and $|y|>R$,
let $a\in F$, and put $j_K=\lfloor KN/L\rfloor$. In the window $z/L\to1$, so
for large $N$ we have $j_K\le N/2$ and $uj_K\ge K/2$. As in Step~2 of the proof
of Lemma~\ref{lem:start-upper}, but with the sum split at $j_K$,
\[
 S^{(a)}_n(u)\le u\sum_{j\le j_K,\,j<n_a}(1-u)^{j-1}S_{n-je_a}(u)
 +\bigl(e^{-K/2}+O(L^{-1})\bigr)\bar S_{N,k}(u).
\]
Now let $j\le j_K$ and $j<n_a$. Put $n'=n-je_a$, $N'=N-j$, $L'=\log N'$ and
$y'=\sqrt{L'}\,(n'/N'-\frac1k\mathbf1)$. The chain of length $N'$ with the same
$p$ has $T'=N'r$, the same $u$, and the window coordinate
$t'=T'-L'+\frac12\log L'$. Since $T-T'=Tj/N\le2K$ and $0\le L-L'\le2j/N\le2K/L$,
the value $t'$ lies in a compact set $\mathcal J'$ that depends only on
$\mathcal J$ and $K$. Each coordinate of $n'/N'-n/N$ has absolute value at most
$j/N'\le2j/N$, so $|n'/N'-n/N|\le2\sqrt k\,K/L$. With $|y|\le2\sqrt L$ and
$\sqrt{L'/L}\ge1-2K/L^2$ this gives $|y'|\ge|y|-O(L^{-1/2})>R/2$ for large $N$.
So Lemma~\ref{lem:confinement}(b) at length $N'$, with $\mathcal J'$ and $R/2$,
and Theorem~\ref{thm:simplex-window}(iii) at length $N'$ give
$S_{n'}(u)\le e^{-k^2R^2/16+o(1)}\bar S_{N',k}(u)$, uniformly. With
\eqref{eq:first-run-sum} we get
\[
 kS^{(a)}_n(u)\le\bigl(e^{-k^2R^2/16}+ke^{-K/2}+o(1)\bigr)\bar S_{N,k}(u).
\]
Choose $K$ with $ke^{-K/2}<\frac12(1-e^{-k^2R^2/16})$. Then, as in (1), there
is $\delta>0$ such that for large $N$ every bin with support $F$ and $|y|>R$ has
$P^\mu_N(n)\le(1-\delta)\frac{\mu(F)}kp^{N-1}\bar S_{N,k}(u)$, while a balanced
bin has $1-o(1)$ times the same value. So no maximizer has $|y|>R$.
\end{proof}

\begin{proof}[Proof of Corollary~\ref{cor:start-modes}]
By Proposition~\ref{prop:dpmf}(b), a vertex $i$ has
$P^\mu_N(Ne_i)=\mu_ip^{N-1}$, and every bin on a face $S$ with $\mu(S)=0$ has
probability $0$. For a face $S$ with $|S|=k\ge2$ and $\mu(S)>0$,
Theorem~\ref{thm:start-maxima}(1) gives the best bin of $S$ the value
$\log\frac{\mu(S)}k+(k-1)(t-a_k)+o(1)$ of $\log(P^\mu_N/p^{N-1})$, uniformly
for $t\in\mathcal J$. Among the faces of size $k$ the largest $\mu(S)$ is
$M_k$. So this value is $\ell^\mu_k(t)+o(1)$ when $\mu(S)=M_k$, and it is at
most $\ell^\mu_k(t)-\delta_1+o(1)$ when $0<\mu(S)<M_k$, where $\delta_1>0$
depends only on $\mu$. The same holds for the vertices, with $\ell^\mu_1$ and
without the term $o(1)$. Since $\mathcal J$ is compact and the lines are
continuous, there is $\delta_2>0$ with
$\ell^\mu_{k^*}(t)\ge\ell^\mu_k(t)+\delta_2$ for all $t\in\mathcal J$ and
$k\ne k^*$. There are finitely many faces. So for large $N$ and all
$t\in\mathcal J$, the best bin of a face of size $k^*$ with $\mu(S)=M_{k^*}$ is
more likely than every bin on every face that does not have size $k^*$ and mass
$\mu(S)=M_{k^*}$. This proves the first claim and
the claim for $k^*=1$. If $\mu$ is constant on $S$, then
$P^\mu_N(n)=\mu_ap^{N-1}S_n(u)$ for $a\in S$ and every $n$ with support $S$, and
$S_n(u)$ is symmetric in the coordinates. Since $p>1/d$ for large $N$,
Theorem~\ref{thm:simplex-balance} gives the claim on balanced bins.
The last claim is Theorem~\ref{thm:start-maxima}(2), since a global mode
maximizes $P^\mu_N$ among the bins with its support, and each coordinate of $y$
is at most $|y|$ in absolute value.
\end{proof}

\subsection{The field}

\begin{proof}[Proof of Theorem~\ref{thm:weakfield-phases}]
Fix compact sets of $t$ and $h$, and a support $S$ with $|S|=k\ge2$.

\emph{Step 1. The best bin on $S$.} Relative to $V_N$, a bin with support $S$
has weight $e^{h\cdot\alpha}S_n(u)$ with $\alpha=n/N$. Choose $R$ with
$k^2R^2/4>\max_ih_i-\bar h_S+1$ on the compact set. Since
$h\cdot\alpha\le\max_{i\in S}h_i$, Lemma~\ref{lem:confinement}(b) shows that
bins with $|y|>R$ have log weight at most $\ell_S(t,h)-1+o(1)$. On $|y|\le R$ we
have $h\cdot\alpha=\bar h_S+O(R/\sqrt L)$, and Lemma~\ref{lem:confinement}(a)
gives the log weight $\ell_S(t,h)-\frac{k^2}4|y|^2+o(1)$. A balanced bin has
$|y|=O(\sqrt L/N)$. So the largest weight on $S$ is $e^{\ell_S(t,h)+o(1)}$,
while a vertex $i$ has weight exactly $e^{h_i}$.

\emph{Step 2. The modes.} There are finitely many supports. This
proves~\eqref{eq:weakfield-maximum}, and a positive gap to a lower score excludes
that support for large $N$. If the maximizers on a selected support had
$|y|\ge\eta>0$ along a subsequence, their log weight would be at most
$\ell_S-k^2\eta^2/4+o(1)$, below that of a balanced bin. This proves the
location statement.

\emph{Step 3. The lines.} Among supports of size $k$ the largest $\bar h_S$ is
the average of the $k$ largest fields, which gives~\eqref{eq:weakfield-lines}. If
$k_1$ wins at $t_1$ and $k_2$ at $t_2>t_1$, adding the two comparisons gives
$(k_1-k_2)(t_2-t_1)\le0$. So the winning size is nondecreasing away from ties.

\emph{Step 4. Every intermediate size.} To realize a given $2\le k<d$, set
$h_1=\cdots=h_k=0$ and $h_{k+1}=\cdots=h_d=-H$, and take $t>a_k$ so close to
$a_k$ that $t<a_{k-1}$ if $k>2$. Then the $k$-face score is positive, the
largest vertex score is $0$, and every face with $2\le j<k$ states scores at
most $(j-1)(t-a_j)<0$. A face with $j>k$ states scores at most
$(j-1)(t-a_j)-H(j-k)/j$, and another face with $k$ states at most
$(k-1)(t-a_k)-H/k$, so both lose once $H$ is large. The comparisons are strict,
so they hold on an open interval of $t$, and small changes make the fields
strictly ordered without removing the phase.
\end{proof}

\begin{proof}[Proof of Corollary~\ref{cor:weakfield-full-hierarchy}]
Since $\frac1k\sum_{i=1}^k(i-1)^2=(k-1)(2k-1)/6$, the lines
\eqref{eq:weakfield-lines} are $\ell_k(t)=(k-1)t+c_k-H(k-1)(2k-1)/6$, also for
$k=1$, and $\ell_{k+1}(t)-\ell_k(t)=t-t_k$. Extend $c_k$ to the smooth function
$c(x)=(x+\frac12)\log x-\frac{x-1}2\log(4\pi)$ for $x\ge1$. Then
$0<c''(x)=(x-\frac12)/x^2\le\frac12$, and integrating twice over unit intervals
gives $c_{k+2}-2c_{k+1}+c_k\le\frac12$. Hence
\[
 t_{k+1}-t_k=\frac{2H}{3}-(c_{k+2}-2c_{k+1}+c_k)>0.
\]
For $t_{k-1}<t<t_k$, summing the differences $\ell_{j+1}(t)-\ell_j(t)=t-t_j$
shows that line $k$ strictly exceeds every other line. The fields are strictly
decreasing, so the unique best support of size $k$ is $\{1,\ldots,k\}$, and
Theorem~\ref{thm:weakfield-phases} gives the conclusion for finite $N$.
\end{proof}
